\documentclass[12pt,a4paper]{article}
\usepackage[top=1.8cm,bottom=1.8cm,left=1.8cm,right=1.8cm]{geometry}
\usepackage{fontspec}
\setmainfont{TeX Gyre Termes}
\usepackage{amsmath,amssymb}
\usepackage{fancyhdr}
\usepackage{tasks}
\usepackage{tikz}
\settasks{label=\textbf{\Alph*.},label-width=1.5em,item-indent=2em}
\setlength{\parindent}{0pt}
\setlength{\parskip}{3pt}
\pagestyle{fancy}\fancyhf{}\renewcommand{\headrulewidth}{0pt}
\fancyhead[L]{\small Ngân hàng bài tập Toán 9 -- 9H2-C6-B18-D03}
\cfoot{\small Trang \thepage}
\newif\ifloigiai
\loigiaitrue
\newenvironment{cauhoi}[2]{%
  \par\medskip\noindent\rule{\textwidth}{0.4pt}\par\nopagebreak
  \noindent\textbf{[#1]}\hfill{\scriptsize\ttfamily #2}\par\nopagebreak\smallskip}{\par}
\newcommand{\loigiai}[1]{\ifloigiai\par\smallskip\noindent\textit{Lời giải.} #1\par\fi}
\newcommand{\dapan}[1]{\ifloigiai\par\smallskip\noindent\textit{Đáp án:} \textbf{#1}\par\fi}
\begin{document}
\begin{center}
{\Large\textbf{NGÂN HÀNG BÀI TẬP TOÁN 9}}\\[3pt]
\textbf{Chương 6. Hàm số y = ax² (a ≠ 0). Phương trình bậc hai một ẩn}\\[2pt]
\textbf{Bài 18. Hàm số y = ax² (a ≠ 0)}\\[2pt]
\textbf{Dạng D03: Tìm m để cắt tại hai điểm thoả điều kiện}\\[2pt]
\textit{Mã dạng: 9H2-C6-B18-D03 \quad 50 bài tự luận + 50 câu trắc nghiệm \quad (VD: 25, VDC: 25)}
\end{center}
\textbf{Quy ước mã câu:} \textit{khối, học kì, chương, bài, dạng, loại (TL/TN), mức độ, số thứ tự}.\par Ví dụ: \texttt{9H2-C6-B18-D03-TL-NB-001}.\par Để in bản cho học sinh, đặt \verb|\loigiaifalse| ngay sau \verb|\loigiaitrue|.
\section*{Phần I. Tự luận}
\begin{cauhoi}{Bài 1}{9H2-C6-B18-D03-TL-VD-001}
Tìm $m$ để đường thẳng $(d): y=-12x+m$ cắt parabol $(P): y=x^{2}$ tại hai điểm phân biệt.
\loigiai{Phương trình hoành độ giao điểm $x^{2}+12x-m=0$. Để $(d)$ cắt $(P)$ tại hai điểm phân biệt cần $\Delta=144+4m>0$, suy ra $m>-36$.}
\end{cauhoi}
\begin{cauhoi}{Bài 2}{9H2-C6-B18-D03-TL-VD-002}
Tìm $m$ để đường thẳng $(d): y=-33x+m$ cắt parabol $(P): y=x^{2}$ tại hai điểm phân biệt.
\loigiai{Phương trình hoành độ giao điểm $x^{2}+33x-m=0$. Để $(d)$ cắt $(P)$ tại hai điểm phân biệt cần $\Delta=1089+4m>0$, suy ra $m>-\dfrac{1089}{4}$.}
\end{cauhoi}
\begin{cauhoi}{Bài 3}{9H2-C6-B18-D03-TL-VD-003}
Tìm $m$ để đường thẳng $(d): y=27x+m$ cắt parabol $(P): y=x^{2}$ tại hai điểm phân biệt.
\loigiai{Phương trình hoành độ giao điểm $x^{2}-27x-m=0$. Để $(d)$ cắt $(P)$ tại hai điểm phân biệt cần $\Delta=729+4m>0$, suy ra $m>-\dfrac{729}{4}$.}
\end{cauhoi}
\begin{cauhoi}{Bài 4}{9H2-C6-B18-D03-TL-VD-004}
Tìm $m$ để đường thẳng $(d): y=24x+m$ cắt parabol $(P): y=x^{2}$ tại hai điểm phân biệt.
\loigiai{Phương trình hoành độ giao điểm $x^{2}-24x-m=0$. Để $(d)$ cắt $(P)$ tại hai điểm phân biệt cần $\Delta=576+4m>0$, suy ra $m>-144$.}
\end{cauhoi}
\begin{cauhoi}{Bài 5}{9H2-C6-B18-D03-TL-VD-005}
Tìm $m$ để đường thẳng $(d): y=23x+m$ cắt parabol $(P): y=x^{2}$ tại hai điểm phân biệt.
\loigiai{Phương trình hoành độ giao điểm $x^{2}-23x-m=0$. Để $(d)$ cắt $(P)$ tại hai điểm phân biệt cần $\Delta=529+4m>0$, suy ra $m>-\dfrac{529}{4}$.}
\end{cauhoi}
\begin{cauhoi}{Bài 6}{9H2-C6-B18-D03-TL-VD-006}
Tìm $m$ để đường thẳng $(d): y=33x+m$ cắt parabol $(P): y=x^{2}$ tại hai điểm phân biệt.
\loigiai{Phương trình hoành độ giao điểm $x^{2}-33x-m=0$. Để $(d)$ cắt $(P)$ tại hai điểm phân biệt cần $\Delta=1089+4m>0$, suy ra $m>-\dfrac{1089}{4}$.}
\end{cauhoi}
\begin{cauhoi}{Bài 7}{9H2-C6-B18-D03-TL-VD-007}
Tìm $m$ để đường thẳng $(d): y=22x+m$ cắt parabol $(P): y=x^{2}$ tại hai điểm phân biệt.
\loigiai{Phương trình hoành độ giao điểm $x^{2}-22x-m=0$. Để $(d)$ cắt $(P)$ tại hai điểm phân biệt cần $\Delta=484+4m>0$, suy ra $m>-121$.}
\end{cauhoi}
\begin{cauhoi}{Bài 8}{9H2-C6-B18-D03-TL-VD-008}
Tìm $m$ để đường thẳng $(d): y=-3x+m$ cắt parabol $(P): y=x^{2}$ tại hai điểm phân biệt.
\loigiai{Phương trình hoành độ giao điểm $x^{2}+3x-m=0$. Để $(d)$ cắt $(P)$ tại hai điểm phân biệt cần $\Delta=9+4m>0$, suy ra $m>-\dfrac{9}{4}$.}
\end{cauhoi}
\begin{cauhoi}{Bài 9}{9H2-C6-B18-D03-TL-VD-009}
Tìm $m$ để đường thẳng $(d): y=-13x+m$ cắt parabol $(P): y=x^{2}$ tại hai điểm phân biệt.
\loigiai{Phương trình hoành độ giao điểm $x^{2}+13x-m=0$. Để $(d)$ cắt $(P)$ tại hai điểm phân biệt cần $\Delta=169+4m>0$, suy ra $m>-\dfrac{169}{4}$.}
\end{cauhoi}
\begin{cauhoi}{Bài 10}{9H2-C6-B18-D03-TL-VD-010}
Tìm $m$ để đường thẳng $(d): y=20x+m$ cắt parabol $(P): y=x^{2}$ tại hai điểm phân biệt.
\loigiai{Phương trình hoành độ giao điểm $x^{2}-20x-m=0$. Để $(d)$ cắt $(P)$ tại hai điểm phân biệt cần $\Delta=400+4m>0$, suy ra $m>-100$.}
\end{cauhoi}
\begin{cauhoi}{Bài 11}{9H2-C6-B18-D03-TL-VD-011}
Tìm $m$ để đường thẳng $(d): y=40x+m$ cắt parabol $(P): y=x^{2}$ tại hai điểm phân biệt.
\loigiai{Phương trình hoành độ giao điểm $x^{2}-40x-m=0$. Để $(d)$ cắt $(P)$ tại hai điểm phân biệt cần $\Delta=1600+4m>0$, suy ra $m>-400$.}
\end{cauhoi}
\begin{cauhoi}{Bài 12}{9H2-C6-B18-D03-TL-VD-012}
Tìm $m$ để đường thẳng $(d): y=-21x+m$ cắt parabol $(P): y=x^{2}$ tại hai điểm phân biệt.
\loigiai{Phương trình hoành độ giao điểm $x^{2}+21x-m=0$. Để $(d)$ cắt $(P)$ tại hai điểm phân biệt cần $\Delta=441+4m>0$, suy ra $m>-\dfrac{441}{4}$.}
\end{cauhoi}
\begin{cauhoi}{Bài 13}{9H2-C6-B18-D03-TL-VD-013}
Tìm $m$ để đường thẳng $(d): y=-20x+m$ cắt parabol $(P): y=x^{2}$ tại hai điểm phân biệt.
\loigiai{Phương trình hoành độ giao điểm $x^{2}+20x-m=0$. Để $(d)$ cắt $(P)$ tại hai điểm phân biệt cần $\Delta=400+4m>0$, suy ra $m>-100$.}
\end{cauhoi}
\begin{cauhoi}{Bài 14}{9H2-C6-B18-D03-TL-VD-014}
Tìm $m$ để đường thẳng $(d): y=34x+m$ cắt parabol $(P): y=x^{2}$ tại hai điểm phân biệt.
\loigiai{Phương trình hoành độ giao điểm $x^{2}-34x-m=0$. Để $(d)$ cắt $(P)$ tại hai điểm phân biệt cần $\Delta=1156+4m>0$, suy ra $m>-289$.}
\end{cauhoi}
\begin{cauhoi}{Bài 15}{9H2-C6-B18-D03-TL-VD-015}
Tìm $m$ để đường thẳng $(d): y=17x+m$ cắt parabol $(P): y=x^{2}$ tại hai điểm phân biệt.
\loigiai{Phương trình hoành độ giao điểm $x^{2}-17x-m=0$. Để $(d)$ cắt $(P)$ tại hai điểm phân biệt cần $\Delta=289+4m>0$, suy ra $m>-\dfrac{289}{4}$.}
\end{cauhoi}
\begin{cauhoi}{Bài 16}{9H2-C6-B18-D03-TL-VD-016}
Tìm $m$ để đường thẳng $(d): y=-23x+m$ cắt parabol $(P): y=x^{2}$ tại hai điểm phân biệt.
\loigiai{Phương trình hoành độ giao điểm $x^{2}+23x-m=0$. Để $(d)$ cắt $(P)$ tại hai điểm phân biệt cần $\Delta=529+4m>0$, suy ra $m>-\dfrac{529}{4}$.}
\end{cauhoi}
\begin{cauhoi}{Bài 17}{9H2-C6-B18-D03-TL-VD-017}
Tìm $m$ để đường thẳng $(d): y=-38x+m$ cắt parabol $(P): y=x^{2}$ tại hai điểm phân biệt.
\loigiai{Phương trình hoành độ giao điểm $x^{2}+38x-m=0$. Để $(d)$ cắt $(P)$ tại hai điểm phân biệt cần $\Delta=1444+4m>0$, suy ra $m>-361$.}
\end{cauhoi}
\begin{cauhoi}{Bài 18}{9H2-C6-B18-D03-TL-VD-018}
Tìm $m$ để đường thẳng $(d): y=18x+m$ cắt parabol $(P): y=x^{2}$ tại hai điểm phân biệt.
\loigiai{Phương trình hoành độ giao điểm $x^{2}-18x-m=0$. Để $(d)$ cắt $(P)$ tại hai điểm phân biệt cần $\Delta=324+4m>0$, suy ra $m>-81$.}
\end{cauhoi}
\begin{cauhoi}{Bài 19}{9H2-C6-B18-D03-TL-VD-019}
Tìm $m$ để đường thẳng $(d): y=15x+m$ cắt parabol $(P): y=x^{2}$ tại hai điểm phân biệt.
\loigiai{Phương trình hoành độ giao điểm $x^{2}-15x-m=0$. Để $(d)$ cắt $(P)$ tại hai điểm phân biệt cần $\Delta=225+4m>0$, suy ra $m>-\dfrac{225}{4}$.}
\end{cauhoi}
\begin{cauhoi}{Bài 20}{9H2-C6-B18-D03-TL-VD-020}
Tìm $m$ để đường thẳng $(d): y=-7x+m$ cắt parabol $(P): y=x^{2}$ tại hai điểm phân biệt.
\loigiai{Phương trình hoành độ giao điểm $x^{2}+7x-m=0$. Để $(d)$ cắt $(P)$ tại hai điểm phân biệt cần $\Delta=49+4m>0$, suy ra $m>-\dfrac{49}{4}$.}
\end{cauhoi}
\begin{cauhoi}{Bài 21}{9H2-C6-B18-D03-TL-VD-021}
Tìm $m$ để đường thẳng $(d): y=21x+m$ cắt parabol $(P): y=x^{2}$ tại hai điểm phân biệt.
\loigiai{Phương trình hoành độ giao điểm $x^{2}-21x-m=0$. Để $(d)$ cắt $(P)$ tại hai điểm phân biệt cần $\Delta=441+4m>0$, suy ra $m>-\dfrac{441}{4}$.}
\end{cauhoi}
\begin{cauhoi}{Bài 22}{9H2-C6-B18-D03-TL-VD-022}
Tìm $m$ để đường thẳng $(d): y=-6x+m$ cắt parabol $(P): y=x^{2}$ tại hai điểm phân biệt.
\loigiai{Phương trình hoành độ giao điểm $x^{2}+6x-m=0$. Để $(d)$ cắt $(P)$ tại hai điểm phân biệt cần $\Delta=36+4m>0$, suy ra $m>-9$.}
\end{cauhoi}
\begin{cauhoi}{Bài 23}{9H2-C6-B18-D03-TL-VD-023}
Tìm $m$ để đường thẳng $(d): y=35x+m$ cắt parabol $(P): y=x^{2}$ tại hai điểm phân biệt.
\loigiai{Phương trình hoành độ giao điểm $x^{2}-35x-m=0$. Để $(d)$ cắt $(P)$ tại hai điểm phân biệt cần $\Delta=1225+4m>0$, suy ra $m>-\dfrac{1225}{4}$.}
\end{cauhoi}
\begin{cauhoi}{Bài 24}{9H2-C6-B18-D03-TL-VD-024}
Tìm $m$ để đường thẳng $(d): y=-15x+m$ cắt parabol $(P): y=x^{2}$ tại hai điểm phân biệt.
\loigiai{Phương trình hoành độ giao điểm $x^{2}+15x-m=0$. Để $(d)$ cắt $(P)$ tại hai điểm phân biệt cần $\Delta=225+4m>0$, suy ra $m>-\dfrac{225}{4}$.}
\end{cauhoi}
\begin{cauhoi}{Bài 25}{9H2-C6-B18-D03-TL-VD-025}
Tìm $m$ để đường thẳng $(d): y=2x+m$ cắt parabol $(P): y=x^{2}$ tại hai điểm phân biệt.
\loigiai{Phương trình hoành độ giao điểm $x^{2}-2x-m=0$. Để $(d)$ cắt $(P)$ tại hai điểm phân biệt cần $\Delta=4+4m>0$, suy ra $m>-1$.}
\end{cauhoi}
\begin{cauhoi}{Bài 26}{9H2-C6-B18-D03-TL-VDC-026}
Tìm $m$ để đường thẳng $(d): y=32x+m$ cắt parabol $(P): y=x^{2}$ tại hai điểm phân biệt.
\loigiai{Phương trình hoành độ giao điểm $x^{2}-32x-m=0$. Để $(d)$ cắt $(P)$ tại hai điểm phân biệt cần $\Delta=1024+4m>0$, suy ra $m>-256$.}
\end{cauhoi}
\begin{cauhoi}{Bài 27}{9H2-C6-B18-D03-TL-VDC-027}
Tìm $m$ để đường thẳng $(d): y=-18x+m$ cắt parabol $(P): y=x^{2}$ tại hai điểm phân biệt.
\loigiai{Phương trình hoành độ giao điểm $x^{2}+18x-m=0$. Để $(d)$ cắt $(P)$ tại hai điểm phân biệt cần $\Delta=324+4m>0$, suy ra $m>-81$.}
\end{cauhoi}
\begin{cauhoi}{Bài 28}{9H2-C6-B18-D03-TL-VDC-028}
Tìm $m$ để đường thẳng $(d): y=19x+m$ cắt parabol $(P): y=x^{2}$ tại hai điểm phân biệt.
\loigiai{Phương trình hoành độ giao điểm $x^{2}-19x-m=0$. Để $(d)$ cắt $(P)$ tại hai điểm phân biệt cần $\Delta=361+4m>0$, suy ra $m>-\dfrac{361}{4}$.}
\end{cauhoi}
\begin{cauhoi}{Bài 29}{9H2-C6-B18-D03-TL-VDC-029}
Tìm $m$ để đường thẳng $(d): y=8x+m$ cắt parabol $(P): y=x^{2}$ tại hai điểm phân biệt.
\loigiai{Phương trình hoành độ giao điểm $x^{2}-8x-m=0$. Để $(d)$ cắt $(P)$ tại hai điểm phân biệt cần $\Delta=64+4m>0$, suy ra $m>-16$.}
\end{cauhoi}
\begin{cauhoi}{Bài 30}{9H2-C6-B18-D03-TL-VDC-030}
Tìm $m$ để đường thẳng $(d): y=-16x+m$ cắt parabol $(P): y=x^{2}$ tại hai điểm phân biệt.
\loigiai{Phương trình hoành độ giao điểm $x^{2}+16x-m=0$. Để $(d)$ cắt $(P)$ tại hai điểm phân biệt cần $\Delta=256+4m>0$, suy ra $m>-64$.}
\end{cauhoi}
\begin{cauhoi}{Bài 31}{9H2-C6-B18-D03-TL-VDC-031}
Tìm $m$ để đường thẳng $(d): y=3x+m$ cắt parabol $(P): y=x^{2}$ tại hai điểm phân biệt.
\loigiai{Phương trình hoành độ giao điểm $x^{2}-3x-m=0$. Để $(d)$ cắt $(P)$ tại hai điểm phân biệt cần $\Delta=9+4m>0$, suy ra $m>-\dfrac{9}{4}$.}
\end{cauhoi}
\begin{cauhoi}{Bài 32}{9H2-C6-B18-D03-TL-VDC-032}
Tìm $m$ để đường thẳng $(d): y=-39x+m$ cắt parabol $(P): y=x^{2}$ tại hai điểm phân biệt.
\loigiai{Phương trình hoành độ giao điểm $x^{2}+39x-m=0$. Để $(d)$ cắt $(P)$ tại hai điểm phân biệt cần $\Delta=1521+4m>0$, suy ra $m>-\dfrac{1521}{4}$.}
\end{cauhoi}
\begin{cauhoi}{Bài 33}{9H2-C6-B18-D03-TL-VDC-033}
Tìm $m$ để đường thẳng $(d): y=31x+m$ cắt parabol $(P): y=x^{2}$ tại hai điểm phân biệt.
\loigiai{Phương trình hoành độ giao điểm $x^{2}-31x-m=0$. Để $(d)$ cắt $(P)$ tại hai điểm phân biệt cần $\Delta=961+4m>0$, suy ra $m>-\dfrac{961}{4}$.}
\end{cauhoi}
\begin{cauhoi}{Bài 34}{9H2-C6-B18-D03-TL-VDC-034}
Tìm $m$ để đường thẳng $(d): y=-9x+m$ cắt parabol $(P): y=x^{2}$ tại hai điểm phân biệt.
\loigiai{Phương trình hoành độ giao điểm $x^{2}+9x-m=0$. Để $(d)$ cắt $(P)$ tại hai điểm phân biệt cần $\Delta=81+4m>0$, suy ra $m>-\dfrac{81}{4}$.}
\end{cauhoi}
\begin{cauhoi}{Bài 35}{9H2-C6-B18-D03-TL-VDC-035}
Tìm $m$ để đường thẳng $(d): y=7x+m$ cắt parabol $(P): y=x^{2}$ tại hai điểm phân biệt.
\loigiai{Phương trình hoành độ giao điểm $x^{2}-7x-m=0$. Để $(d)$ cắt $(P)$ tại hai điểm phân biệt cần $\Delta=49+4m>0$, suy ra $m>-\dfrac{49}{4}$.}
\end{cauhoi}
\begin{cauhoi}{Bài 36}{9H2-C6-B18-D03-TL-VDC-036}
Tìm $m$ để đường thẳng $(d): y=38x+m$ cắt parabol $(P): y=x^{2}$ tại hai điểm phân biệt.
\loigiai{Phương trình hoành độ giao điểm $x^{2}-38x-m=0$. Để $(d)$ cắt $(P)$ tại hai điểm phân biệt cần $\Delta=1444+4m>0$, suy ra $m>-361$.}
\end{cauhoi}
\begin{cauhoi}{Bài 37}{9H2-C6-B18-D03-TL-VDC-037}
Tìm $m$ để đường thẳng $(d): y=9x+m$ cắt parabol $(P): y=x^{2}$ tại hai điểm phân biệt.
\loigiai{Phương trình hoành độ giao điểm $x^{2}-9x-m=0$. Để $(d)$ cắt $(P)$ tại hai điểm phân biệt cần $\Delta=81+4m>0$, suy ra $m>-\dfrac{81}{4}$.}
\end{cauhoi}
\begin{cauhoi}{Bài 38}{9H2-C6-B18-D03-TL-VDC-038}
Tìm $m$ để đường thẳng $(d): y=39x+m$ cắt parabol $(P): y=x^{2}$ tại hai điểm phân biệt.
\loigiai{Phương trình hoành độ giao điểm $x^{2}-39x-m=0$. Để $(d)$ cắt $(P)$ tại hai điểm phân biệt cần $\Delta=1521+4m>0$, suy ra $m>-\dfrac{1521}{4}$.}
\end{cauhoi}
\begin{cauhoi}{Bài 39}{9H2-C6-B18-D03-TL-VDC-039}
Tìm $m$ để đường thẳng $(d): y=-29x+m$ cắt parabol $(P): y=x^{2}$ tại hai điểm phân biệt.
\loigiai{Phương trình hoành độ giao điểm $x^{2}+29x-m=0$. Để $(d)$ cắt $(P)$ tại hai điểm phân biệt cần $\Delta=841+4m>0$, suy ra $m>-\dfrac{841}{4}$.}
\end{cauhoi}
\begin{cauhoi}{Bài 40}{9H2-C6-B18-D03-TL-VDC-040}
Tìm $m$ để đường thẳng $(d): y=-x+m$ cắt parabol $(P): y=x^{2}$ tại hai điểm phân biệt.
\loigiai{Phương trình hoành độ giao điểm $x^{2}+x-m=0$. Để $(d)$ cắt $(P)$ tại hai điểm phân biệt cần $\Delta=1+4m>0$, suy ra $m>-\dfrac{1}{4}$.}
\end{cauhoi}
\begin{cauhoi}{Bài 41}{9H2-C6-B18-D03-TL-VDC-041}
Tìm $m$ để đường thẳng $(d): y=x+m$ cắt parabol $(P): y=x^{2}$ tại hai điểm phân biệt.
\loigiai{Phương trình hoành độ giao điểm $x^{2}-x-m=0$. Để $(d)$ cắt $(P)$ tại hai điểm phân biệt cần $\Delta=1+4m>0$, suy ra $m>-\dfrac{1}{4}$.}
\end{cauhoi}
\begin{cauhoi}{Bài 42}{9H2-C6-B18-D03-TL-VDC-042}
Tìm $m$ để đường thẳng $(d): y=-10x+m$ cắt parabol $(P): y=x^{2}$ tại hai điểm phân biệt.
\loigiai{Phương trình hoành độ giao điểm $x^{2}+10x-m=0$. Để $(d)$ cắt $(P)$ tại hai điểm phân biệt cần $\Delta=100+4m>0$, suy ra $m>-25$.}
\end{cauhoi}
\begin{cauhoi}{Bài 43}{9H2-C6-B18-D03-TL-VDC-043}
Tìm $m$ để đường thẳng $(d): y=-40x+m$ cắt parabol $(P): y=x^{2}$ tại hai điểm phân biệt.
\loigiai{Phương trình hoành độ giao điểm $x^{2}+40x-m=0$. Để $(d)$ cắt $(P)$ tại hai điểm phân biệt cần $\Delta=1600+4m>0$, suy ra $m>-400$.}
\end{cauhoi}
\begin{cauhoi}{Bài 44}{9H2-C6-B18-D03-TL-VDC-044}
Tìm $m$ để đường thẳng $(d): y=-30x+m$ cắt parabol $(P): y=x^{2}$ tại hai điểm phân biệt.
\loigiai{Phương trình hoành độ giao điểm $x^{2}+30x-m=0$. Để $(d)$ cắt $(P)$ tại hai điểm phân biệt cần $\Delta=900+4m>0$, suy ra $m>-225$.}
\end{cauhoi}
\begin{cauhoi}{Bài 45}{9H2-C6-B18-D03-TL-VDC-045}
Tìm $m$ để đường thẳng $(d): y=-36x+m$ cắt parabol $(P): y=x^{2}$ tại hai điểm phân biệt.
\loigiai{Phương trình hoành độ giao điểm $x^{2}+36x-m=0$. Để $(d)$ cắt $(P)$ tại hai điểm phân biệt cần $\Delta=1296+4m>0$, suy ra $m>-324$.}
\end{cauhoi}
\begin{cauhoi}{Bài 46}{9H2-C6-B18-D03-TL-VDC-046}
Tìm $m$ để đường thẳng $(d): y=13x+m$ cắt parabol $(P): y=x^{2}$ tại hai điểm phân biệt.
\loigiai{Phương trình hoành độ giao điểm $x^{2}-13x-m=0$. Để $(d)$ cắt $(P)$ tại hai điểm phân biệt cần $\Delta=169+4m>0$, suy ra $m>-\dfrac{169}{4}$.}
\end{cauhoi}
\begin{cauhoi}{Bài 47}{9H2-C6-B18-D03-TL-VDC-047}
Tìm $m$ để đường thẳng $(d): y=30x+m$ cắt parabol $(P): y=x^{2}$ tại hai điểm phân biệt.
\loigiai{Phương trình hoành độ giao điểm $x^{2}-30x-m=0$. Để $(d)$ cắt $(P)$ tại hai điểm phân biệt cần $\Delta=900+4m>0$, suy ra $m>-225$.}
\end{cauhoi}
\begin{cauhoi}{Bài 48}{9H2-C6-B18-D03-TL-VDC-048}
Tìm $m$ để đường thẳng $(d): y=-19x+m$ cắt parabol $(P): y=x^{2}$ tại hai điểm phân biệt.
\loigiai{Phương trình hoành độ giao điểm $x^{2}+19x-m=0$. Để $(d)$ cắt $(P)$ tại hai điểm phân biệt cần $\Delta=361+4m>0$, suy ra $m>-\dfrac{361}{4}$.}
\end{cauhoi}
\begin{cauhoi}{Bài 49}{9H2-C6-B18-D03-TL-VDC-049}
Tìm $m$ để đường thẳng $(d): y=11x+m$ cắt parabol $(P): y=x^{2}$ tại hai điểm phân biệt.
\loigiai{Phương trình hoành độ giao điểm $x^{2}-11x-m=0$. Để $(d)$ cắt $(P)$ tại hai điểm phân biệt cần $\Delta=121+4m>0$, suy ra $m>-\dfrac{121}{4}$.}
\end{cauhoi}
\begin{cauhoi}{Bài 50}{9H2-C6-B18-D03-TL-VDC-050}
Tìm $m$ để đường thẳng $(d): y=6x+m$ cắt parabol $(P): y=x^{2}$ tại hai điểm phân biệt.
\loigiai{Phương trình hoành độ giao điểm $x^{2}-6x-m=0$. Để $(d)$ cắt $(P)$ tại hai điểm phân biệt cần $\Delta=36+4m>0$, suy ra $m>-9$.}
\end{cauhoi}
\clearpage\section*{Phần II. Trắc nghiệm}
\begin{cauhoi}{Câu 1}{9H2-C6-B18-D03-TN-VD-001}
Tìm $m$ để đường thẳng $(d): y=-12x+m$ cắt parabol $(P): y=x^{2}$ tại hai điểm phân biệt.
\begin{tasks}(2)
\task $m>-36$
\task $m\ge -36$
\task $m>36$
\task $m<-36$
\end{tasks}
\dapan{A}
\end{cauhoi}
\begin{cauhoi}{Câu 2}{9H2-C6-B18-D03-TN-VD-002}
Tìm $m$ để đường thẳng $(d): y=-33x+m$ cắt parabol $(P): y=x^{2}$ tại hai điểm phân biệt.
\begin{tasks}(2)
\task $m>\dfrac{1089}{4}$
\task $m\ge -\dfrac{1089}{4}$
\task $m<-\dfrac{1089}{4}$
\task $m>-\dfrac{1089}{4}$
\end{tasks}
\dapan{D}
\end{cauhoi}
\begin{cauhoi}{Câu 3}{9H2-C6-B18-D03-TN-VD-003}
Tìm $m$ để đường thẳng $(d): y=27x+m$ cắt parabol $(P): y=x^{2}$ tại hai điểm phân biệt.
\begin{tasks}(2)
\task $m>\dfrac{729}{4}$
\task $m>-\dfrac{729}{4}$
\task $m\ge -\dfrac{729}{4}$
\task $m<-\dfrac{729}{4}$
\end{tasks}
\dapan{B}
\end{cauhoi}
\begin{cauhoi}{Câu 4}{9H2-C6-B18-D03-TN-VD-004}
Tìm $m$ để đường thẳng $(d): y=24x+m$ cắt parabol $(P): y=x^{2}$ tại hai điểm phân biệt.
\begin{tasks}(2)
\task $m<-144$
\task $m>144$
\task $m>-144$
\task $m\ge -144$
\end{tasks}
\dapan{C}
\end{cauhoi}
\begin{cauhoi}{Câu 5}{9H2-C6-B18-D03-TN-VD-005}
Tìm $m$ để đường thẳng $(d): y=23x+m$ cắt parabol $(P): y=x^{2}$ tại hai điểm phân biệt.
\begin{tasks}(2)
\task $m<-\dfrac{529}{4}$
\task $m>-\dfrac{529}{4}$
\task $m>\dfrac{529}{4}$
\task $m\ge -\dfrac{529}{4}$
\end{tasks}
\dapan{B}
\end{cauhoi}
\begin{cauhoi}{Câu 6}{9H2-C6-B18-D03-TN-VD-006}
Tìm $m$ để đường thẳng $(d): y=33x+m$ cắt parabol $(P): y=x^{2}$ tại hai điểm phân biệt.
\begin{tasks}(2)
\task $m>\dfrac{1089}{4}$
\task $m<-\dfrac{1089}{4}$
\task $m\ge -\dfrac{1089}{4}$
\task $m>-\dfrac{1089}{4}$
\end{tasks}
\dapan{D}
\end{cauhoi}
\begin{cauhoi}{Câu 7}{9H2-C6-B18-D03-TN-VD-007}
Tìm $m$ để đường thẳng $(d): y=22x+m$ cắt parabol $(P): y=x^{2}$ tại hai điểm phân biệt.
\begin{tasks}(2)
\task $m\ge -121$
\task $m>-121$
\task $m<-121$
\task $m>121$
\end{tasks}
\dapan{B}
\end{cauhoi}
\begin{cauhoi}{Câu 8}{9H2-C6-B18-D03-TN-VD-008}
Tìm $m$ để đường thẳng $(d): y=-3x+m$ cắt parabol $(P): y=x^{2}$ tại hai điểm phân biệt.
\begin{tasks}(2)
\task $m\ge -\dfrac{9}{4}$
\task $m<-\dfrac{9}{4}$
\task $m>-\dfrac{9}{4}$
\task $m>\dfrac{9}{4}$
\end{tasks}
\dapan{C}
\end{cauhoi}
\begin{cauhoi}{Câu 9}{9H2-C6-B18-D03-TN-VD-009}
Tìm $m$ để đường thẳng $(d): y=-13x+m$ cắt parabol $(P): y=x^{2}$ tại hai điểm phân biệt.
\begin{tasks}(2)
\task $m>-\dfrac{169}{4}$
\task $m\ge -\dfrac{169}{4}$
\task $m<-\dfrac{169}{4}$
\task $m>\dfrac{169}{4}$
\end{tasks}
\dapan{A}
\end{cauhoi}
\begin{cauhoi}{Câu 10}{9H2-C6-B18-D03-TN-VD-010}
Tìm $m$ để đường thẳng $(d): y=20x+m$ cắt parabol $(P): y=x^{2}$ tại hai điểm phân biệt.
\begin{tasks}(2)
\task $m>100$
\task $m\ge -100$
\task $m<-100$
\task $m>-100$
\end{tasks}
\dapan{D}
\end{cauhoi}
\begin{cauhoi}{Câu 11}{9H2-C6-B18-D03-TN-VD-011}
Tìm $m$ để đường thẳng $(d): y=40x+m$ cắt parabol $(P): y=x^{2}$ tại hai điểm phân biệt.
\begin{tasks}(2)
\task $m<-400$
\task $m>400$
\task $m>-400$
\task $m\ge -400$
\end{tasks}
\dapan{C}
\end{cauhoi}
\begin{cauhoi}{Câu 12}{9H2-C6-B18-D03-TN-VD-012}
Tìm $m$ để đường thẳng $(d): y=-21x+m$ cắt parabol $(P): y=x^{2}$ tại hai điểm phân biệt.
\begin{tasks}(2)
\task $m>\dfrac{441}{4}$
\task $m>-\dfrac{441}{4}$
\task $m\ge -\dfrac{441}{4}$
\task $m<-\dfrac{441}{4}$
\end{tasks}
\dapan{B}
\end{cauhoi}
\begin{cauhoi}{Câu 13}{9H2-C6-B18-D03-TN-VD-013}
Tìm $m$ để đường thẳng $(d): y=-20x+m$ cắt parabol $(P): y=x^{2}$ tại hai điểm phân biệt.
\begin{tasks}(2)
\task $m<-100$
\task $m>100$
\task $m>-100$
\task $m\ge -100$
\end{tasks}
\dapan{C}
\end{cauhoi}
\begin{cauhoi}{Câu 14}{9H2-C6-B18-D03-TN-VD-014}
Tìm $m$ để đường thẳng $(d): y=34x+m$ cắt parabol $(P): y=x^{2}$ tại hai điểm phân biệt.
\begin{tasks}(2)
\task $m>-289$
\task $m<-289$
\task $m>289$
\task $m\ge -289$
\end{tasks}
\dapan{A}
\end{cauhoi}
\begin{cauhoi}{Câu 15}{9H2-C6-B18-D03-TN-VD-015}
Tìm $m$ để đường thẳng $(d): y=17x+m$ cắt parabol $(P): y=x^{2}$ tại hai điểm phân biệt.
\begin{tasks}(2)
\task $m>-\dfrac{289}{4}$
\task $m>\dfrac{289}{4}$
\task $m<-\dfrac{289}{4}$
\task $m\ge -\dfrac{289}{4}$
\end{tasks}
\dapan{A}
\end{cauhoi}
\begin{cauhoi}{Câu 16}{9H2-C6-B18-D03-TN-VD-016}
Tìm $m$ để đường thẳng $(d): y=-23x+m$ cắt parabol $(P): y=x^{2}$ tại hai điểm phân biệt.
\begin{tasks}(2)
\task $m>-\dfrac{529}{4}$
\task $m\ge -\dfrac{529}{4}$
\task $m>\dfrac{529}{4}$
\task $m<-\dfrac{529}{4}$
\end{tasks}
\dapan{A}
\end{cauhoi}
\begin{cauhoi}{Câu 17}{9H2-C6-B18-D03-TN-VD-017}
Tìm $m$ để đường thẳng $(d): y=-38x+m$ cắt parabol $(P): y=x^{2}$ tại hai điểm phân biệt.
\begin{tasks}(2)
\task $m\ge -361$
\task $m>361$
\task $m<-361$
\task $m>-361$
\end{tasks}
\dapan{D}
\end{cauhoi}
\begin{cauhoi}{Câu 18}{9H2-C6-B18-D03-TN-VD-018}
Tìm $m$ để đường thẳng $(d): y=18x+m$ cắt parabol $(P): y=x^{2}$ tại hai điểm phân biệt.
\begin{tasks}(2)
\task $m>81$
\task $m<-81$
\task $m>-81$
\task $m\ge -81$
\end{tasks}
\dapan{C}
\end{cauhoi}
\begin{cauhoi}{Câu 19}{9H2-C6-B18-D03-TN-VD-019}
Tìm $m$ để đường thẳng $(d): y=15x+m$ cắt parabol $(P): y=x^{2}$ tại hai điểm phân biệt.
\begin{tasks}(2)
\task $m>-\dfrac{225}{4}$
\task $m>\dfrac{225}{4}$
\task $m\ge -\dfrac{225}{4}$
\task $m<-\dfrac{225}{4}$
\end{tasks}
\dapan{A}
\end{cauhoi}
\begin{cauhoi}{Câu 20}{9H2-C6-B18-D03-TN-VD-020}
Tìm $m$ để đường thẳng $(d): y=-7x+m$ cắt parabol $(P): y=x^{2}$ tại hai điểm phân biệt.
\begin{tasks}(2)
\task $m\ge -\dfrac{49}{4}$
\task $m>-\dfrac{49}{4}$
\task $m>\dfrac{49}{4}$
\task $m<-\dfrac{49}{4}$
\end{tasks}
\dapan{B}
\end{cauhoi}
\begin{cauhoi}{Câu 21}{9H2-C6-B18-D03-TN-VD-021}
Tìm $m$ để đường thẳng $(d): y=21x+m$ cắt parabol $(P): y=x^{2}$ tại hai điểm phân biệt.
\begin{tasks}(2)
\task $m<-\dfrac{441}{4}$
\task $m\ge -\dfrac{441}{4}$
\task $m>-\dfrac{441}{4}$
\task $m>\dfrac{441}{4}$
\end{tasks}
\dapan{C}
\end{cauhoi}
\begin{cauhoi}{Câu 22}{9H2-C6-B18-D03-TN-VD-022}
Tìm $m$ để đường thẳng $(d): y=-6x+m$ cắt parabol $(P): y=x^{2}$ tại hai điểm phân biệt.
\begin{tasks}(2)
\task $m<-9$
\task $m\ge -9$
\task $m>9$
\task $m>-9$
\end{tasks}
\dapan{D}
\end{cauhoi}
\begin{cauhoi}{Câu 23}{9H2-C6-B18-D03-TN-VD-023}
Tìm $m$ để đường thẳng $(d): y=35x+m$ cắt parabol $(P): y=x^{2}$ tại hai điểm phân biệt.
\begin{tasks}(2)
\task $m<-\dfrac{1225}{4}$
\task $m\ge -\dfrac{1225}{4}$
\task $m>\dfrac{1225}{4}$
\task $m>-\dfrac{1225}{4}$
\end{tasks}
\dapan{D}
\end{cauhoi}
\begin{cauhoi}{Câu 24}{9H2-C6-B18-D03-TN-VD-024}
Tìm $m$ để đường thẳng $(d): y=-15x+m$ cắt parabol $(P): y=x^{2}$ tại hai điểm phân biệt.
\begin{tasks}(2)
\task $m>\dfrac{225}{4}$
\task $m>-\dfrac{225}{4}$
\task $m<-\dfrac{225}{4}$
\task $m\ge -\dfrac{225}{4}$
\end{tasks}
\dapan{B}
\end{cauhoi}
\begin{cauhoi}{Câu 25}{9H2-C6-B18-D03-TN-VD-025}
Tìm $m$ để đường thẳng $(d): y=2x+m$ cắt parabol $(P): y=x^{2}$ tại hai điểm phân biệt.
\begin{tasks}(2)
\task $m>-1$
\task $m<-1$
\task $m\ge -1$
\task $m>1$
\end{tasks}
\dapan{A}
\end{cauhoi}
\begin{cauhoi}{Câu 26}{9H2-C6-B18-D03-TN-VDC-026}
Tìm $m$ để đường thẳng $(d): y=32x+m$ cắt parabol $(P): y=x^{2}$ tại hai điểm phân biệt.
\begin{tasks}(2)
\task $m<-256$
\task $m\ge -256$
\task $m>256$
\task $m>-256$
\end{tasks}
\dapan{D}
\end{cauhoi}
\begin{cauhoi}{Câu 27}{9H2-C6-B18-D03-TN-VDC-027}
Tìm $m$ để đường thẳng $(d): y=-18x+m$ cắt parabol $(P): y=x^{2}$ tại hai điểm phân biệt.
\begin{tasks}(2)
\task $m>81$
\task $m\ge -81$
\task $m>-81$
\task $m<-81$
\end{tasks}
\dapan{C}
\end{cauhoi}
\begin{cauhoi}{Câu 28}{9H2-C6-B18-D03-TN-VDC-028}
Tìm $m$ để đường thẳng $(d): y=19x+m$ cắt parabol $(P): y=x^{2}$ tại hai điểm phân biệt.
\begin{tasks}(2)
\task $m<-\dfrac{361}{4}$
\task $m\ge -\dfrac{361}{4}$
\task $m>-\dfrac{361}{4}$
\task $m>\dfrac{361}{4}$
\end{tasks}
\dapan{C}
\end{cauhoi}
\begin{cauhoi}{Câu 29}{9H2-C6-B18-D03-TN-VDC-029}
Tìm $m$ để đường thẳng $(d): y=8x+m$ cắt parabol $(P): y=x^{2}$ tại hai điểm phân biệt.
\begin{tasks}(2)
\task $m<-16$
\task $m>16$
\task $m\ge -16$
\task $m>-16$
\end{tasks}
\dapan{D}
\end{cauhoi}
\begin{cauhoi}{Câu 30}{9H2-C6-B18-D03-TN-VDC-030}
Tìm $m$ để đường thẳng $(d): y=-16x+m$ cắt parabol $(P): y=x^{2}$ tại hai điểm phân biệt.
\begin{tasks}(2)
\task $m>64$
\task $m\ge -64$
\task $m<-64$
\task $m>-64$
\end{tasks}
\dapan{D}
\end{cauhoi}
\begin{cauhoi}{Câu 31}{9H2-C6-B18-D03-TN-VDC-031}
Tìm $m$ để đường thẳng $(d): y=3x+m$ cắt parabol $(P): y=x^{2}$ tại hai điểm phân biệt.
\begin{tasks}(2)
\task $m>\dfrac{9}{4}$
\task $m\ge -\dfrac{9}{4}$
\task $m<-\dfrac{9}{4}$
\task $m>-\dfrac{9}{4}$
\end{tasks}
\dapan{D}
\end{cauhoi}
\begin{cauhoi}{Câu 32}{9H2-C6-B18-D03-TN-VDC-032}
Tìm $m$ để đường thẳng $(d): y=-39x+m$ cắt parabol $(P): y=x^{2}$ tại hai điểm phân biệt.
\begin{tasks}(2)
\task $m>-\dfrac{1521}{4}$
\task $m<-\dfrac{1521}{4}$
\task $m>\dfrac{1521}{4}$
\task $m\ge -\dfrac{1521}{4}$
\end{tasks}
\dapan{A}
\end{cauhoi}
\begin{cauhoi}{Câu 33}{9H2-C6-B18-D03-TN-VDC-033}
Tìm $m$ để đường thẳng $(d): y=31x+m$ cắt parabol $(P): y=x^{2}$ tại hai điểm phân biệt.
\begin{tasks}(2)
\task $m>-\dfrac{961}{4}$
\task $m>\dfrac{961}{4}$
\task $m\ge -\dfrac{961}{4}$
\task $m<-\dfrac{961}{4}$
\end{tasks}
\dapan{A}
\end{cauhoi}
\begin{cauhoi}{Câu 34}{9H2-C6-B18-D03-TN-VDC-034}
Tìm $m$ để đường thẳng $(d): y=-9x+m$ cắt parabol $(P): y=x^{2}$ tại hai điểm phân biệt.
\begin{tasks}(2)
\task $m\ge -\dfrac{81}{4}$
\task $m<-\dfrac{81}{4}$
\task $m>-\dfrac{81}{4}$
\task $m>\dfrac{81}{4}$
\end{tasks}
\dapan{C}
\end{cauhoi}
\begin{cauhoi}{Câu 35}{9H2-C6-B18-D03-TN-VDC-035}
Tìm $m$ để đường thẳng $(d): y=7x+m$ cắt parabol $(P): y=x^{2}$ tại hai điểm phân biệt.
\begin{tasks}(2)
\task $m<-\dfrac{49}{4}$
\task $m>\dfrac{49}{4}$
\task $m\ge -\dfrac{49}{4}$
\task $m>-\dfrac{49}{4}$
\end{tasks}
\dapan{D}
\end{cauhoi}
\begin{cauhoi}{Câu 36}{9H2-C6-B18-D03-TN-VDC-036}
Tìm $m$ để đường thẳng $(d): y=38x+m$ cắt parabol $(P): y=x^{2}$ tại hai điểm phân biệt.
\begin{tasks}(2)
\task $m>-361$
\task $m>361$
\task $m<-361$
\task $m\ge -361$
\end{tasks}
\dapan{A}
\end{cauhoi}
\begin{cauhoi}{Câu 37}{9H2-C6-B18-D03-TN-VDC-037}
Tìm $m$ để đường thẳng $(d): y=9x+m$ cắt parabol $(P): y=x^{2}$ tại hai điểm phân biệt.
\begin{tasks}(2)
\task $m<-\dfrac{81}{4}$
\task $m\ge -\dfrac{81}{4}$
\task $m>-\dfrac{81}{4}$
\task $m>\dfrac{81}{4}$
\end{tasks}
\dapan{C}
\end{cauhoi}
\begin{cauhoi}{Câu 38}{9H2-C6-B18-D03-TN-VDC-038}
Tìm $m$ để đường thẳng $(d): y=39x+m$ cắt parabol $(P): y=x^{2}$ tại hai điểm phân biệt.
\begin{tasks}(2)
\task $m>-\dfrac{1521}{4}$
\task $m\ge -\dfrac{1521}{4}$
\task $m>\dfrac{1521}{4}$
\task $m<-\dfrac{1521}{4}$
\end{tasks}
\dapan{A}
\end{cauhoi}
\begin{cauhoi}{Câu 39}{9H2-C6-B18-D03-TN-VDC-039}
Tìm $m$ để đường thẳng $(d): y=-29x+m$ cắt parabol $(P): y=x^{2}$ tại hai điểm phân biệt.
\begin{tasks}(2)
\task $m>\dfrac{841}{4}$
\task $m>-\dfrac{841}{4}$
\task $m<-\dfrac{841}{4}$
\task $m\ge -\dfrac{841}{4}$
\end{tasks}
\dapan{B}
\end{cauhoi}
\begin{cauhoi}{Câu 40}{9H2-C6-B18-D03-TN-VDC-040}
Tìm $m$ để đường thẳng $(d): y=-x+m$ cắt parabol $(P): y=x^{2}$ tại hai điểm phân biệt.
\begin{tasks}(2)
\task $m>\dfrac{1}{4}$
\task $m>-\dfrac{1}{4}$
\task $m\ge -\dfrac{1}{4}$
\task $m<-\dfrac{1}{4}$
\end{tasks}
\dapan{B}
\end{cauhoi}
\begin{cauhoi}{Câu 41}{9H2-C6-B18-D03-TN-VDC-041}
Tìm $m$ để đường thẳng $(d): y=x+m$ cắt parabol $(P): y=x^{2}$ tại hai điểm phân biệt.
\begin{tasks}(2)
\task $m>-\dfrac{1}{4}$
\task $m<-\dfrac{1}{4}$
\task $m\ge -\dfrac{1}{4}$
\task $m>\dfrac{1}{4}$
\end{tasks}
\dapan{A}
\end{cauhoi}
\begin{cauhoi}{Câu 42}{9H2-C6-B18-D03-TN-VDC-042}
Tìm $m$ để đường thẳng $(d): y=-10x+m$ cắt parabol $(P): y=x^{2}$ tại hai điểm phân biệt.
\begin{tasks}(2)
\task $m>-25$
\task $m>25$
\task $m<-25$
\task $m\ge -25$
\end{tasks}
\dapan{A}
\end{cauhoi}
\begin{cauhoi}{Câu 43}{9H2-C6-B18-D03-TN-VDC-043}
Tìm $m$ để đường thẳng $(d): y=-40x+m$ cắt parabol $(P): y=x^{2}$ tại hai điểm phân biệt.
\begin{tasks}(2)
\task $m>-400$
\task $m>400$
\task $m<-400$
\task $m\ge -400$
\end{tasks}
\dapan{A}
\end{cauhoi}
\begin{cauhoi}{Câu 44}{9H2-C6-B18-D03-TN-VDC-044}
Tìm $m$ để đường thẳng $(d): y=-30x+m$ cắt parabol $(P): y=x^{2}$ tại hai điểm phân biệt.
\begin{tasks}(2)
\task $m<-225$
\task $m>225$
\task $m>-225$
\task $m\ge -225$
\end{tasks}
\dapan{C}
\end{cauhoi}
\begin{cauhoi}{Câu 45}{9H2-C6-B18-D03-TN-VDC-045}
Tìm $m$ để đường thẳng $(d): y=-36x+m$ cắt parabol $(P): y=x^{2}$ tại hai điểm phân biệt.
\begin{tasks}(2)
\task $m\ge -324$
\task $m>-324$
\task $m>324$
\task $m<-324$
\end{tasks}
\dapan{B}
\end{cauhoi}
\begin{cauhoi}{Câu 46}{9H2-C6-B18-D03-TN-VDC-046}
Tìm $m$ để đường thẳng $(d): y=13x+m$ cắt parabol $(P): y=x^{2}$ tại hai điểm phân biệt.
\begin{tasks}(2)
\task $m<-\dfrac{169}{4}$
\task $m>-\dfrac{169}{4}$
\task $m>\dfrac{169}{4}$
\task $m\ge -\dfrac{169}{4}$
\end{tasks}
\dapan{B}
\end{cauhoi}
\begin{cauhoi}{Câu 47}{9H2-C6-B18-D03-TN-VDC-047}
Tìm $m$ để đường thẳng $(d): y=30x+m$ cắt parabol $(P): y=x^{2}$ tại hai điểm phân biệt.
\begin{tasks}(2)
\task $m>225$
\task $m\ge -225$
\task $m>-225$
\task $m<-225$
\end{tasks}
\dapan{C}
\end{cauhoi}
\begin{cauhoi}{Câu 48}{9H2-C6-B18-D03-TN-VDC-048}
Tìm $m$ để đường thẳng $(d): y=-19x+m$ cắt parabol $(P): y=x^{2}$ tại hai điểm phân biệt.
\begin{tasks}(2)
\task $m<-\dfrac{361}{4}$
\task $m>\dfrac{361}{4}$
\task $m\ge -\dfrac{361}{4}$
\task $m>-\dfrac{361}{4}$
\end{tasks}
\dapan{D}
\end{cauhoi}
\begin{cauhoi}{Câu 49}{9H2-C6-B18-D03-TN-VDC-049}
Tìm $m$ để đường thẳng $(d): y=11x+m$ cắt parabol $(P): y=x^{2}$ tại hai điểm phân biệt.
\begin{tasks}(2)
\task $m<-\dfrac{121}{4}$
\task $m\ge -\dfrac{121}{4}$
\task $m>\dfrac{121}{4}$
\task $m>-\dfrac{121}{4}$
\end{tasks}
\dapan{D}
\end{cauhoi}
\begin{cauhoi}{Câu 50}{9H2-C6-B18-D03-TN-VDC-050}
Tìm $m$ để đường thẳng $(d): y=6x+m$ cắt parabol $(P): y=x^{2}$ tại hai điểm phân biệt.
\begin{tasks}(2)
\task $m\ge -9$
\task $m<-9$
\task $m>-9$
\task $m>9$
\end{tasks}
\dapan{C}
\end{cauhoi}
\clearpage\section*{Bảng đáp án trắc nghiệm}
\begin{center}\begin{tabular}{|c|c|c|c|c|c|c|c|c|c|}
\hline
1. A & 2. D & 3. B & 4. C & 5. B & 6. D & 7. B & 8. C & 9. A & 10. D\\\hline
11. C & 12. B & 13. C & 14. A & 15. A & 16. A & 17. D & 18. C & 19. A & 20. B\\\hline
21. C & 22. D & 23. D & 24. B & 25. A & 26. D & 27. C & 28. C & 29. D & 30. D\\\hline
31. D & 32. A & 33. A & 34. C & 35. D & 36. A & 37. C & 38. A & 39. B & 40. B\\\hline
41. A & 42. A & 43. A & 44. C & 45. B & 46. B & 47. C & 48. D & 49. D & 50. C\\\hline
\end{tabular}\end{center}
\end{document}
